% Vietnamese translation for OI Public Library
% Source-PDF: 其他 OTHERS/CDQ分治入门经典.pdf
\documentclass[11pt]{article}
\usepackage{oipl-vn}

\newcommand{\Oh}{\mathcal{O}}
\newcommand{\midp}{\mathrm{mid}}
\newcommand{\Max}{\operatorname{Max}}
\newcommand{\Min}{\operatorname{Min}}
\newcommand{\Solve}{\operatorname{Solve}}
\newcommand{\SolveTwo}{\operatorname{Solve2}}
\newcommand{\slope}{\operatorname{slope}}
\newcommand{\dis}{\operatorname{dis}}

\title{Nhập môn kinh điển về phân trị CDQ}
\author{AC\_Aerolight}
\date{28--29 tháng 4, 2013 @ Fuzhou}

\begin{document}
\maketitle

\origtitle{Discussion topic: divide-and-conquer methods}
\sourcepath{Others / CDQ divide-and-conquer introduction.pdf}

\begin{abstract}
Bài giảng trình bày một kiểu phân trị đặc biệt, thường được cộng đồng OI gọi
là phân trị CDQ. Tư tưởng chính là chia miền chỉ số thành hai nửa, giải nửa
trái trước, xử lý toàn bộ ảnh hưởng của nửa trái lên nửa phải, rồi tiếp tục
giải nửa phải. Cách nhìn này biến nhiều bài toán trực tuyến khó xử lý thành
các bài toán offline có thể sắp xếp, quét và dùng cấu trúc dữ liệu chuẩn.
Nguồn là bộ slide dạng ảnh; bản tiếng Việt này chuyển các slide thành ghi chú
bài giảng, giữ mạch ví dụ, các công thức chuyển và các truy hồi độ phức tạp
chính.
\end{abstract}

\tableofcontents

\section{Lời nói đầu}

Tài liệu gốc nhắc tới một dạng phân trị đặc biệt. Cách làm này được thấy
rộng rãi từ bài tập đội tuyển quốc gia năm 2008 của Chen Danqi, nên trong
giới OI Trung Quốc nó thường được gọi là phân trị CDQ.

Mục đích của bài giảng là phổ biến phương pháp: các ví dụ được chọn tương
đối cơ bản, phần suy luận được viết chậm rãi, và những bài đã quen thuộc có
thể được đọc lướt để tập trung vào mẫu tư duy chung.

Một vài công thức truy hồi xuất hiện nhiều lần:
\[
  T(n)=2T(n/2)+\Oh(kn)\Rightarrow T(n)=\Oh(kn\log n),
\]
\[
  T(n)=2T(n/2)+\Oh(kn\log n)\Rightarrow
  T(n)=\Oh(kn\log^2 n),
\]
\[
  T(n)=2T(n/2)+\Oh(k)\Rightarrow T(n)=\Oh(kn).
\]
Trong đó \(k\) là một đa thức không phụ thuộc vào \(n\).

\section{Ví dụ mở đầu: đếm nghịch thế}

Với một hoán vị \(P\), cần đếm số cặp nghịch thế
\[
  (i,j)\quad\text{thỏa}\quad i<j,\ P_i>P_j.
\]
Nếu \(|P|\le 100000\), thuật toán \(\Oh(n^2)\) là quá chậm. Merge sort cho
cách làm \(\Oh(n\log n)\):
\begin{enumerate}
  \item Gọi đệ quy trên nửa trái.
  \item Gọi đệ quy trên nửa phải.
  \item Đếm các cặp có chỉ số trái nằm ở nửa trái và chỉ số phải nằm ở nửa
  phải trong lúc trộn hai dãy đã sắp xếp.
\end{enumerate}
Điểm quan trọng là ảnh hưởng nội bộ của từng nửa đã được đệ quy giải quyết;
bước gộp chỉ cần quan tâm tới ảnh hưởng chéo.

\section{NOI2007 Cash}

\subsection{Quy hoạch động ban đầu}

Có hai loại trái phiếu. Khi mua ở ngày \(i\), tiền được chia theo tỉ lệ
\(\operatorname{Rate}[i]\) giữa hai loại; khi bán, người chơi bán một tỉ lệ
cố định của lượng trái phiếu đang giữ. Biết trước giá hai loại trái phiếu
\(A[i]\), \(B[i]\) trong \(n\) ngày và vốn ban đầu \(S\), cần tối đa hóa lợi
nhuận. Để tối ưu, mỗi lần giao dịch rõ ràng nên mua hết hoặc bán hết. Ràng
buộc chính là
\[
  1\le n\le 100000.
\]

Đặt \(F[i]\) là số lượng trái phiếu loại \(B\) lớn nhất có thể có sau ngày
\(i\). Khi xét giao dịch trước đó là ngày \(j\), công thức trong slide có
dạng
\[
  F[i]=\max\left\{F[i],
  \frac{\operatorname{Rate}[j]F[j]A[i]+F[j]B[i]}
       {\operatorname{Rate}[i]A[i]+B[i]}\right\}.
\]
So sánh hai quyết định \(j,k\), ta nhận được điều kiện
\[
  \frac{\operatorname{Rate}[j]F[j]-\operatorname{Rate}[k]F[k]}
       {F[j]-F[k]}
  > -\frac{B[i]}{A[i]}.
\]
Vế trái thường được gọi là \(\slope(j,k)\).

Đặt
\[
  G[i]=\operatorname{Rate}[i]F[i],\qquad X_i=(F[i],G[i]).
\]
Khi đó \(\slope(j,k)\) là hệ số góc của đường thẳng qua hai điểm \(X_j,X_k\).
Bài toán cần duy trì một tập điểm động hỗ trợ:
\begin{enumerate}
  \item chèn một điểm bất kỳ trong góc phần tư thứ nhất;
  \item với một hệ số góc âm \(K\), tìm đường thẳng có hệ số góc \(K\), đi
  qua một điểm của tập, và có tung độ cắt trục \(Y\) lớn nhất.
\end{enumerate}
Những điểm có thể được chọn ở thao tác thứ hai đều nằm trên bao lồi trên của
tập điểm. Nếu xử lý trực tuyến, số biên nhiều và cài đặt dễ rối.

\subsection{Đưa phân trị vào}

Quan sát then chốt là hệ số góc truy vấn \(-B[i]/A[i]\) chỉ phụ thuộc vào
ngày \(i\), không phụ thuộc vào giá trị hiện tại của \(F\). Khi chèn điểm
\(X_i\), ta đã biết nó sẽ ảnh hưởng tới những truy vấn \(j>i\) nào. Vì vậy
có thể định nghĩa
\[
  \Solve(l,r)
\]
là thủ tục tính được các giá trị \(F[l],F[l+1],\ldots,F[r]\).

Với \(\midp=(l+r)/2\), thủ tục có dạng:
\begin{enumerate}
  \item chạy \(\Solve(l,\midp)\);
  \item xét đồng thời toàn bộ ảnh hưởng của các điểm trong \([l,\midp]\)
  lên các truy vấn trong \([\midp+1,r]\);
  \item chạy \(\Solve(\midp+1,r)\).
\end{enumerate}
Như vậy, phần giữa chỉ cần bao lồi của điểm ở nửa trái và các truy vấn ở nửa
phải. Bao lồi của hai tập đã sắp xếp có thể gộp tuyến tính; mỗi truy vấn trên
bao lồi được trả lời bằng nhị phân. Bước giữa tốn \(\Oh(n\log n)\) trên đoạn
độ dài \(n=r-l+1\), do đó
\[
  T(n)=2T(n/2)+\Oh(n\log n)=\Oh(n\log^2 n).
\]

Có thể tiếp tục giảm bước giữa xuống tuyến tính. Thay vì trả lời từng truy
vấn bằng nhị phân trên bao lồi, ta bỏ nhị phân và xử lý offline: sắp xếp bao
lồi điểm cùng toàn bộ truy vấn theo hệ số góc, rồi dùng hai con trỏ quét để
tìm điểm tối ưu cho từng truy vấn. Muốn làm được điều đó trong mọi nút đệ quy,
cần sắp xếp trước các truy vấn và lưu kết quả của các lần trộn trung gian; nếu
đợi đến khi gọi \(\Solve\) mới trộn thì vẫn lặp lại chi phí sắp xếp.

Chi phí tiền xử lý trộn là \(\Oh(n\log n)\). Trong đệ quy chính, mỗi bước giữa
chỉ còn \(\Oh(n)\), nên
\[
  T(n)=2T(n/2)+\Oh(n)=\Oh(n\log n).
\]
Nếu muốn khử đệ quy, slide gợi ý mô phỏng bằng ngăn xếp thủ công hoặc lưu lại
các quá trình trung gian. Bài giảng kết luận hai ý: phân trị chỉ xét ảnh hưởng
chéo, còn ảnh hưởng nội bộ của mỗi đoạn được đệ quy xử lý; nếu ảnh hưởng nội
bộ có thể tạm bỏ qua, bài toán trở thành offline.

\section{WF2011 Machine Works}

Công ty có quyền sử dụng nhà máy trong \(N\) ngày và vốn ban đầu \(C\). Có
\(M\) máy, mỗi máy có bốn tham số \(D,P,R,G\). Ta có thể mua máy ở ngày
\(D\) với giá \(P\), từ ngày \(D+1\) máy sinh lời \(G\) mỗi ngày, và khi
không cần nữa có thể bán máy thu \(R\). Trong nhà máy chỉ được giữ một máy.
Không thể vận hành máy vào đúng ngày mua hoặc bán, nhưng có thể bán một máy
rồi mua máy khác trong cùng ngày. Sau ngày \(N+1\), phải bán máy đang giữ và
cần tối đa hóa vốn.

Dữ liệu có thể tới
\[
  D\le 10^9,\quad C\le 10^9,\quad M\le 10^5,
\]
\[
  D_i\le D,\quad 1\le R_i<P_i\le 10^9,\quad 1\le G_i\le 10^9.
\]

Sắp xếp máy theo ngày bán \(D_i\), rời rạc hóa thời gian, và đặt \(F[i]\) là
số tiền lớn nhất còn lại sau khi bán máy thứ \(i\). Khi chuyển từ máy \(j\)
sang máy \(i\):
\[
  F[i]=\Max\left(F[i-1],
  F[j]-P[j]+R[j]+G[j](D[i]-D[j]-1)\right),
\]
với điều kiện \(F[j]>P[j]\). Đặt
\[
  A[j]=F[j]-P[j]+R[j]-G[j]D[j]-G[j],
\]
ta có
\[
  F[i]=\Max\left(F[i-1],A[j]+G[j]D[i]\right).
\]
Đây lại là truy vấn cực đại trên một tập đường thẳng
\[
  y=G[j]x+A[j].
\]
Nói theo hình học, cần duy trì một tập đường thẳng với hai thao tác:
\begin{enumerate}
  \item chèn hàm bậc nhất \(y=kx+b\);
  \item với truy vấn \(D\), lấy giá trị lớn nhất của mọi hàm tại \(x=D\).
\end{enumerate}
Slide mô tả đây là việc duy trì một dãy nửa mặt phẳng/bao lồi tuyến tính. Vì
các hoành độ truy vấn \(D[i]\) không đổi theo quá trình tính, ta có thể dùng
phân trị trên toàn bộ đoạn \([1,n]\) giống bài Cash. Mỗi bước giữa duy trì một
họ đường thẳng và trả lời các truy vấn \(x=D[i]\), đạt \(\Oh(n\log n)\), cùng
bậc với cách duy trì bao lồi bằng cây cân bằng.

\section{BOI2007 Mokia}

Bài toán có một bàn cờ rất lớn \(2000000\times 2000000\). Mỗi ô chứa một giá
trị \(A[x,y]\). Hai thao tác:
\[
  \texttt{Add x y a}: A[x,y]\mathrel{+}=a,
\]
\[
  \texttt{Query x0 y0 x1 y1}: \text{hỏi tổng trong hình chữ nhật.}
\]
Số thao tác thêm không quá \(160000\), số truy vấn không quá \(10000\).

Do kích thước bàn cờ và số truy vấn chênh nhau lớn, cấu trúc dữ liệu hai
chiều trực tiếp dễ bị quá bộ nhớ hoặc quá thời gian. Phân trị CDQ xử lý ảnh
hưởng của các thao tác \texttt{Add} ở nửa trái lên các truy vấn ở nửa phải.
Mỗi truy vấn hình chữ nhật được tách bằng hiệu tiền tố hai chiều:
\[
  [x_0,x_1]\times[y_0,y_1]
\]
được đổi thành bốn truy vấn tiền tố \((x,y)\). Ở bước gộp, sắp xếp tất cả
điểm và truy vấn theo \(x\), quét tăng \(x\), và dùng Fenwick hoặc cây đoạn
trên trục \(y\). Khi mỗi lớp phân trị tốn \(\Oh(n\log n)\),
\[
  T(n)=2T(n/2)+\Oh(n\log n)=\Oh(n\log^2 n).
\]

\section{Khoảng cách Manhattan động}

Tài liệu nêu một bài toán duy trì tập điểm hai chiều \(P\), hỗ trợ:
\begin{enumerate}
  \item chèn điểm \((x,y)\);
  \item với truy vấn \((x,y)\), tìm điểm gần nhất theo khoảng cách Manhattan.
\end{enumerate}
Khoảng cách là
\[
  \dis(P_1,P_2)=|x_1-x_2|+|y_1-y_2|,
\]
với \(N,m\le 300000\), \(x,y\le 1000000\). Cây \(k\)-d không phù hợp với
yêu cầu trong bài giảng.

Loại bỏ dấu giá trị tuyệt đối bằng cách xét bốn hướng: trái trên, trái
dưới, phải trên, phải dưới. Chỉ cần mô tả một hướng, ví dụ đáp án nằm ở
phía trái dưới truy vấn. Khi đó
\[
  \dis((x,y),(x',y'))=(x-x')+(y-y')=x+y-(x'+y').
\]
Vì vậy cần tìm điểm thỏa \(x'\le x, y'\le y\) có \(x'+y'\) lớn nhất. Nếu
không có thao tác chèn động, ta chỉ cần sắp xếp theo \(x\) và dùng Fenwick
trên \(y\). Với thao tác chèn theo thời gian, phân trị CDQ một lần nữa xử lý
ảnh hưởng của các điểm ở nửa trái lên truy vấn ở nửa phải:
\begin{enumerate}
  \item \(\Solve(l,\midp)\);
  \item xử lý ảnh hưởng chéo bằng quét offline theo \(x\) và Fenwick theo
  \(y\);
  \item \(\Solve(\midp+1,r)\).
\end{enumerate}
Mỗi hướng được xử lý tương tự, rồi lấy giá trị nhỏ nhất trong bốn đáp án.

\section{Thứ tự từng phần hai và ba chiều}

\subsection{Hai chiều}

Có \(N\) người, mỗi người có hai chỉ số năng lực \(P_i,Q_i\). Nếu
\[
  P_i>P_j\quad\text{và}\quad Q_i>Q_j,
\]
nói rằng \(i\) mạnh hơn \(j\). Cần tìm dãy dài nhất
\[
  A=(A_1,A_2,\ldots,A_t)
\]
sao cho \(A_i\) mạnh hơn \(A_{i+1}\). Với \(N\le 100000\), để dễ nhìn ta giả
sử \(P,Q\) đều là hoán vị của \(1,\ldots,N\).

Sắp xếp các phần tử theo \(P_i\). Bài toán trở thành LIS trên dãy \(Q_i\):
\[
  F[i]=\Max\{F[j]\mid j<i,\ Q[j]<Q[i]\}+1.
\]
Fenwick hoặc cây đoạn cho độ phức tạp \(\Oh(n\log n)\).

\subsection{Ba chiều}

Với ba chỉ số \(P_i,Q_i,R_i\), sau khi sắp xếp theo \(P_i\), cần tính
\[
  F[i]=\Max\{F[j]\mid j<i,\ Q[j]<Q[i],\ R[j]<R[i]\}+1.
\]
Nếu dùng cây đoạn lồng cây cân bằng hoặc cây đoạn bền vững có thể đạt
\(\Oh(n\log^2 n)\). Bài giảng đưa CDQ vào để xử lý chiều chỉ số \(i\):
\[
  \Solve(l,r):
  \begin{cases}
    \Solve(l,\midp),\\
    \text{xử lý ảnh hưởng của }[l,\midp]\text{ lên }[\midp+1,r],\\
    \Solve(\midp+1,r).
  \end{cases}
\]
Ở bước giữa, xem các điểm bên trái là các điểm có trọng số
\[
  X=(Q_i,R_i),\qquad \text{trọng số }F[i],
\]
và với mỗi điểm bên phải \((Q_j,R_j)\), cần tìm trọng số lớn nhất trong các
điểm có
\[
  Q_i<Q_j,\qquad R_i<R_j.
\]
Sắp xếp tất cả điểm và truy vấn theo \(Q\), quét theo \(Q\), và dùng Fenwick
hoặc cây đoạn để lấy cực đại theo \(R\). Do bước giữa tốn \(\Oh(n\log n)\),
\[
  T(n)=2T(n/2)+\Oh(n\log n)=\Oh(n\log^2 n).
\]

\section{Thứ tự từng phần bốn chiều}

Với bốn chỉ số \(P_i,Q_i,R_i,S_i\), có hai biến thể trong slide.

Biến thể thứ nhất chỉ cần với mỗi người \(i\) tìm xem có người nào mạnh hơn:
\[
  j<i,\quad Q_j<Q_i,\quad R_j<R_i,\quad S_j<S_i.
\]
Bài toán tương đương tìm
\[
  \Min\{S_j\mid j<i,\ Q_j<Q_i,\ R_j<R_i\}<S_i.
\]
Sau khi sắp xếp theo \(P\), có thể dùng đúng cách xử lý của bài ba chiều:
CDQ theo \(i\), quét theo \(Q\), và cấu trúc dữ liệu trên \(R\) để giữ giá
trị nhỏ nhất của \(S\).

Biến thể thứ hai yêu cầu dãy dài nhất:
\[
  F[i]=\Max\{F[j]\mid j<i,\ Q_j<Q_i,\ R_j<R_i,\ S_j<S_i\}+1.
\]
Không còn biến đổi đơn giản như trường hợp chỉ hỏi tồn tại, nên bài giảng
dùng hai tầng phân trị. Thủ tục ngoài \(\Solve(l,r)\) giải các giá trị
\(F[l],\ldots,F[r]\). Ở giữa, nó gọi một thủ tục phụ \(\SolveTwo(l,r)\) để
xử lý ảnh hưởng của \([l,\midp]\) lên \([\midp+1,r]\) theo ba chiều còn lại.

Trong \(\SolveTwo\), ta lại phân trị một lần nữa, rồi ở bước giữa của tầng
phụ chỉ còn bài toán hai chiều: chèn điểm \((R_i,S_i)\) với trọng số
\(F[i]\), và truy vấn tại \((R_j,S_j)\) để lấy trọng số lớn nhất trong các
điểm thỏa
\[
  R_i<R_j,\qquad S_i<S_j.
\]
Phần này được xử lý offline bằng cách sắp xếp và dùng Fenwick/cây đoạn giống
trường hợp ba chiều.

Đặt \(T_2(n)\) là độ phức tạp của \(\SolveTwo\):
\[
  T_2(n)=2T_2(n/2)+\Oh(n\log n)=\Oh(n\log^2 n).
\]
Đặt \(T(n)\) là độ phức tạp của \(\Solve\):
\[
  T(n)=2T(n/2)+T_2(n)
      =2T(n/2)+\Oh(n\log^2 n)
      =\Oh(n\log^3 n).
\]

\section{Nhiều chiều và phân trị song song}

Slide đặt câu hỏi vui với thứ tự từng phần \(100\) chiều: mỗi người có
\(100\) chỉ số năng lực \(P_{i,1},P_{i,2},\ldots,P_{i,100}\); nếu mọi chiều
của \(i\) đều lớn hơn của \(j\) thì \(i\) mạnh hơn \(j\). Với \(N\le 500\),
nếu cứ lồng thêm phân trị hoặc cấu trúc dữ liệu theo từng chiều thì sẽ xuất
hiện rất nhiều thừa số \(\log n\). Đây là lời nhắc rằng CDQ là một công cụ
chứ không phải câu trả lời tự động cho mọi số chiều.

Một ý tưởng khác trong bài giảng là phân trị song song. Ví dụ trong bài
POI2011 Meteor có \(N\) quốc gia cùng khai thác quỹ đạo tiểu hành tinh và đặt
\(M\) trạm quan sát. Mỗi trạm thuộc đúng một quốc gia; trạm \(i\) kề trạm
\(i+1\), còn trạm \(M\) kề trạm \(1\). Trong tương lai xảy ra \(K\) trận mưa
sao băng. Trận thứ \(i\) có đoạn \([L_i,R_i]\) theo chiều kim đồng hồ; nếu
\(L_i>R_i\) thì đoạn đi qua cuối vòng, tức gồm \(L_i,L_i+1,\ldots,M,1,\ldots,R_i\).
Mỗi trạm trong đoạn nhận thêm \(A_i\) đơn vị thiên thạch. Quốc gia \(x\) có
mục tiêu \(W_x\); cần in trận mưa đầu tiên sau đó quốc gia này đạt mục tiêu,
hoặc \texttt{NIE} nếu không bao giờ đạt. Slide ghi
\[
  N,M,K\le 3\cdot 10^5,\qquad A_i,W_i\le 10^9.
\]

Ví dụ nguồn:
\begin{verbatim}
3 5
1 3 2 1 3
10 5 7
3
4 2 4
1 3 1
3 5 2
\end{verbatim}
kết quả là
\begin{verbatim}
3
NIE
1
\end{verbatim}

Nếu chỉ có một quốc gia, ta có thể dùng cây đoạn có lazy tag để mô phỏng từng
trận mưa và kiểm tra các trạm của quốc gia đó. Nhị phân đáp án cho một quốc
gia cần mô phỏng \(\Oh(\log K)\) lần. Với nhiều quốc gia, định nghĩa
\(\Solve(l,r,S)\) là thủ tục xử lý tập quốc gia \(S\) mà đáp án nằm trong
\([l,r]\). Sau khi mô phỏng đến \(\midp\), chia \(S\) thành \(S_1\) đã đạt
mục tiêu và \(S_2\) chưa đạt, rồi gọi
\[
  \Solve(l,\midp,S_1),\qquad \Solve(\midp+1,r,S_2).
\]
Cách chia thô phải tính lại quá nhiều lần: phần kiểm tra từng quốc gia tốn
\(\Oh(t_x\log M)\), với \(t_x\) là số trạm của quốc gia \(x\), nhưng phần mô
phỏng các trận mưa có thể dẫn đến truy hồi \(T(x)=2T(x/2)+\Oh(K\log M)\), chưa
đủ tốt.

Cải tiến trong slide là duy trì một cây đoạn toàn cục. Trước khi chạy
\(\Solve(l,r,S)\), cây chứa hiệu ứng của các trận mưa từ \(1\) đến \(l-1\);
sau khi chạy xong, cây chứa hiệu ứng từ \(1\) đến \(r\). Ký hiệu
\(+[l,r]\) là mô phỏng các trận mưa trong đoạn này và \(-[l,r]\) là hoàn tác.
Thủ tục làm:
\[
\begin{array}{l}
\Solve(l,r,S):\\
\quad +[l,\midp],\\
\quad \text{chia }S\text{ theo việc đã đủ mục tiêu ở }\midp,\\
\quad -[l,\midp],\\
\quad \Solve(l,\midp,S_1),\\
\quad \Solve(\midp+1,r,S_2).
\end{array}
\]
Khi đó chi phí mô phỏng theo chiều đáp án là
\[
  T(x)=2T(x/2)+\Oh(x\log M)=\Oh(x\log x\log M),
\]
và toàn bộ phương pháp đạt dạng \(\Oh(K\log K\log M+N\log K\log M)\), đủ qua.

Nhìn theo phân trị song song, ban đầu miền đáp án của mỗi quốc gia là
\([1,K]\). Ta đồng thời nhị phân đáp án của tất cả quốc gia:
\begin{enumerate}
  \item lấy trung điểm \(\operatorname{mid}_x\) của miền đáp án của quốc gia
  \(x\);
  \item sắp xếp các \(\operatorname{mid}_x\), rồi xử lý offline các trận mưa
  sao theo thứ tự thời gian;
  \item nếu tại \(\operatorname{mid}_x\) lượng thu được đủ \(W_x\), đẩy miền
  về nửa trái; ngược lại đẩy về nửa phải.
\end{enumerate}
Sau \(\log k\) vòng, mỗi miền đáp án có độ dài \(1\). Slide ghi độ phức tạp
theo dạng
\[
  \log k\cdot(k\log m+n\log n).
\]

\section{K-th trong ma trận con}

Một ví dụ khác trong slide là bài toán truy vấn thứ tự trên ma trận:
cho ma trận \(A\) kích thước \(n\times n\), có \(q\) truy vấn
\[
  (x_1,y_1,x_2,y_2,k),
\]
hỏi số nhỏ thứ \(k\) trong hình chữ nhật
\([x_1,y_1]\) đến \([x_2,y_2]\). Ràng buộc minh họa là
\[
  n\le 500,\qquad q\le 60000,
\]
và cho phép xử lý offline.

Với một truy vấn đơn lẻ, có thể nhị phân đáp án. Với một giá trị
\(\mathit{mid}\), đặt
\[
  D_{i,j}=[A_{i,j}\le \mathit{mid}],
\]
rồi dùng tổng tiền tố hai chiều để kiểm tra trong hình chữ nhật có ít nhất
\(k\) phần tử không vượt quá \(\mathit{mid}\) hay không.

Để xử lý nhiều truy vấn, định nghĩa \(\Solve(l,r,S)\) là thủ tục xử lý tập
truy vấn \(S\) mà đáp án nằm trong miền giá trị \([l,r]\). Tách tại
\(\mathit{mid}\), dùng cấu trúc dữ liệu để tính với từng truy vấn có bao
nhiêu phần tử \(\le\mathit{mid}\), rồi đưa truy vấn sang nửa trái hoặc nửa
phải. Nếu dùng trực tiếp cây Fenwick hai chiều, mỗi thao tác có thể xem như
\(\Oh(\log^2 n)\), dẫn đến truy hồi dạng
\[
  F(n)=2F(n/2)+\Oh(q\log^2 n).
\]

Slide cũng nhắc cách tối ưu để tránh cấu trúc dữ liệu hai chiều: biến các ô
được chèn và các truy vấn hình chữ nhật thành các sự kiện, sắp xếp/ghép theo
tọa độ \(x\), rồi chỉ cần một cây Fenwick theo tọa độ \(y\). Khi đó bước giữa
còn \(\Oh(n\log n)\) theo số sự kiện của tầng đang xét, và độ phức tạp toàn
cục giảm một thừa số \(\log n\).

\section{ZJOI2013 K-th Query}

Bài toán có \(n\) vị trí và \(m\) thao tác. Có hai loại thao tác:
\begin{itemize}
  \item \(\texttt{1 a b c}\): thêm một số \(c\) vào mọi vị trí từ \(a\) đến
  \(b\);
  \item \(\texttt{2 a b c}\): hỏi trong tất cả các số đang nằm ở các vị trí
  từ \(a\) đến \(b\), số lớn thứ \(c\) là bao nhiêu.
\end{itemize}
Các ràng buộc trong slide là \(n,m\le 50000\); với thao tác thêm thì
\(c\le n\), còn trong truy vấn \(c\) có thể lớn tới miền số nguyên dài.

Cách dùng cấu trúc dữ liệu trực tiếp là cây đoạn lồng cây cân bằng: tầng ngoài
theo giá trị được chèn, tầng trong theo tọa độ vị trí, thường cần cấp phát
động. Với một truy vấn đơn, ta lại có thể nhị phân đáp án và duy trì cây đoạn
hoặc Fenwick để kiểm tra số phần tử trong đoạn.

CDQ theo miền giá trị cho ta một cách làm offline. Gọi \(\Solve(l,r,S)\) xử
lý mọi thao tác/truy vấn trong \(S\) có đáp án thuộc \([l,r]\). Ở giữa
\(\mathit{mid}\), chỉ các thao tác thêm có \(c\le\mathit{mid}\) mới ảnh hưởng
tới câu hỏi ``có bao nhiêu số không vượt quá \(\mathit{mid}\) trong đoạn
\([a,b]\)?''. Ta sắp các thao tác theo thời gian và dùng Fenwick theo vị trí
để cộng/xóa các đoạn tác động, rồi chia \(S\) sang hai nửa.

Sau khi phân bổ chi phí mỗi thao tác ở một tầng, bước giữa ở mức
\(\Oh(n\log n)\), nên truy hồi vẫn có dạng quen thuộc
\[
  T(n)=2T(n/2)+\Oh(n\log n)=\Oh(n\log^2 n).
\]
Ba ví dụ liên tiếp trong phần này có cùng điểm chung: biến câu hỏi thứ tự
thành nhiều phép kiểm tra theo ngưỡng, rồi dùng phân trị trên miền đáp án để
xử lý toàn bộ truy vấn offline.

Slide mẫu cho dữ liệu:
\begin{verbatim}
2 5
1 1 2 1
1 1 2 2
2 1 1 2
2 1 1 1
2 1 2 3
\end{verbatim}
và kết quả:
\begin{verbatim}
1
2
1
\end{verbatim}
Sau thao tác đầu, vị trí \(1\) và \(2\) đều chỉ có số \(1\). Sau thao tác
thứ hai, hai vị trí đó đều có \(1,2\). Truy vấn thứ ba hỏi số lớn thứ \(2\)
trong vị trí \(1\), đáp án \(1\); truy vấn thứ tư hỏi số lớn thứ \(1\) trong
vị trí \(1\), đáp án \(2\); truy vấn thứ năm hỏi số lớn thứ \(3\) trong đoạn
\([1,2]\), đáp án \(1\).

\section{Một mẫu bài khác: Package}

Với bài \emph{Package}, có \(N\) vật, vật \(i\) có trọng lượng \(W_i\) và giá
trị \(G_i\); mỗi vật chỉ được lấy một lần. Có \(Q\) truy vấn, mỗi truy vấn là
\((X,i)\): với ba lô dung lượng tối đa \(X\), dùng tất cả vật trừ vật thứ
\(i\), hỏi tổng giá trị lớn nhất đạt được. Ràng buộc trong slide là
\[
  N\le 100,\qquad X\le 10000,\qquad Q\le 1000000.
\]

Nếu chỉ có một truy vấn, hoặc mọi truy vấn có cùng \(i\), ta chạy quy hoạch
động ba lô \(0/1\) trong \(\Oh(nm)\), với \(m\) là dung lượng. Đưa một vật
vào trạng thái ba lô hiện tại tốn \(\Oh(m)\), nên ký hiệu \(+[l,r]\) là thêm
các vật \(l,\ldots,r\) vào ba lô, chi phí \(\Oh((r-l+1)m)\).

Định nghĩa \(\Solve(l,r)\) để xử lý mọi truy vấn có chỉ số \(l\le i\le r\).
Thủ tục mở rộng có thêm trạng thái \(S\):
\[
  \Solve(l,r,S),
\]
trong đó \(S\) là một ba lô \(0/1\) đã chứa mọi vật nằm ngoài đoạn \([l,r]\).
Khi tới lá \(\Solve(L,L,S)\), ta trả lời tất cả truy vấn \(i=L\).

Đệ quy:
\[
  \Solve(l,r,S):
  \begin{cases}
    \Solve(l,\midp,S+[\midp+1,r]),\\
    \Solve(\midp+1,r,S+[l,\midp]).
  \end{cases}
\]
Nếu mỗi lần gộp trạng thái ba lô tốn \(\Oh(nm)\), truy hồi là
\[
  T(n)=2T(n/2)+\Oh(nm)=\Oh(nm\log n).
\]
Đây là một ví dụ khác của cùng khuôn: tại mỗi nhánh, ta đưa toàn bộ ảnh
hưởng của phần còn lại vào trạng thái, để câu hỏi tại lá trở nên trực tiếp.

\section{HNOI2010 City}

Bài toán cho đồ thị vô hướng \(n\) đỉnh, \(m\) cạnh, mỗi cạnh có trọng số
\(w_i\). Có \(q\) thao tác \((E_i,Z_i)\), nghĩa là đổi trọng số cạnh thứ
\(E_i\) thành \(Z_i\). Sau mỗi thao tác cần in trọng lượng cây khung nhỏ nhất
hiện tại. Ràng buộc trong slide:
\[
  n\le 20000,\quad m\le 50000,\quad q\le 50000,\quad
  w_i,Z_i\le 10^8.
\]

Nếu trọng số của một cạnh chỉ giảm, khi cập nhật cạnh \((x,y)\), cây khung nhỏ
nhất hoặc không đổi, hoặc cạnh \((x,y)\) được thêm vào và một cạnh lớn nhất
trên đường đi từ \(x\) đến \(y\) trong cây cũ bị loại. Khi đó có thể duy trì cây
khung bằng Link-Cut Tree: thêm một cạnh, xóa một cạnh, và hỏi cạnh lớn nhất
trên đường đi. Trường hợp trọng số chỉ tăng là đối xứng hơn nhưng vẫn khó nếu
muốn xử lý trực tuyến đơn giản.

Ở đây lợi thế là toàn bộ dãy thao tác đã biết trước. Ta phân trị trên thời
gian. Gọi \(S\) là tập cạnh bị sửa trong đoạn thao tác đang xét. Trước khi đệ
quy xuống đoạn con, ta muốn thu nhỏ đồ thị mà vẫn giữ đúng mọi đáp án trong
đoạn đó. Slide dùng hai phép: \term{reduction} và \term{contraction}.

\subsection{Reduction}

Cho đồ thị có trọng số \(G\) và một tập cạnh \(S\). Đặt trọng số các cạnh
trong \(S\) thành \(+\infty\), rồi lấy cây khung nhỏ nhất \(T\) của \(G\).
Khi đó dù trọng số các cạnh trong \(S\) thay đổi thế nào, cây khung nhỏ nhất
của \(G\) không cần dùng bất kỳ cạnh nào ngoài \(T\cup S\). Thật vậy, với một
cạnh \(e\notin T\cup S\), hai đầu của \(e\) nối với nhau bằng một đường trong
\(T\). Các cạnh trên đường này đều không thuộc \(S\), nên trọng số của chúng
không đổi khi ta thay đổi \(S\); theo tính chất chu trình của MST, \(e\) không
thể trở thành cạnh cần thiết.

Vì vậy ta có thể xóa mọi cạnh ngoài \(T\cup S\). Số cạnh còn lại không vượt
quá \(n+|S|-1\). Bước này gọi là reduction, mục tiêu là giảm số cạnh. Chi phí
cùng bậc với chạy MST, thường viết \(\Oh(m\log m)\).

\subsection{Contraction}

Ngược lại, đặt trọng số các cạnh trong \(S\) thành \(-\infty\), rồi lấy cây
khung nhỏ nhất \(T\). Khi đó dù trọng số các cạnh trong \(S\) thay đổi thế
nào, mọi cạnh trong \(T\setminus S\) chắc chắn thuộc cây khung nhỏ nhất cuối
cùng. Có thể hình dung tăng dần trọng số từng cạnh của \(S\): mỗi lần thay đổi,
cây khung hoặc không đổi, hoặc một cạnh trong \(S\) bị thay thế; các cạnh
ngoài \(S\) đã được chọn trong \(T\) không bị loại ra theo quá trình này.

Do đó ta có thể co các cạnh trong \(T\setminus S\) thành các thành phần liên
thông và cập nhật lại danh sách cạnh theo thành phần. Bước này gọi là
contraction, mục tiêu là giảm số đỉnh. Sau contraction, số đỉnh còn khoảng
\(|S|+1\).

\subsection{Phân trị theo thao tác}

Từ hai phép trên, với một đồ thị \(G\) và tập thao tác \(S\) của một đoạn thời
gian, ta có thể trong \(\Oh(m\log m)\) rút đồ thị về số cạnh không quá
\(n+|S|-1\) và số đỉnh không quá \(|S|+1\) mà vẫn giữ nguyên kết quả trong
đoạn đang xét.

Định nghĩa \(\Solve(l,r,G)\) để xử lý các thao tác từ \(l\) đến \(r\) trên đồ
thị đã được rút gọn phù hợp với đoạn đó:
\[
\begin{array}{l}
\Solve(l,r,G):\\
\quad \text{reduction}(G);\\
\quad \text{contraction}(G);\\
\quad \Solve(l,\midp,G);\\
\quad \Solve(\midp+1,r,G);\\
\quad \text{recover}(G).
\end{array}
\]
Tại lá, sau một thao tác, đồ thị đã nhỏ tới mức chỉ còn trường hợp cơ bản; có
thể trực tiếp xét hai đỉnh và một cạnh tương ứng để cập nhật đáp án. Trong
mỗi lời gọi \(\Solve(l,r,G)\), sau reduction--contraction, số đỉnh và số cạnh
của \(G\) đều là \(\Oh(r-l)\). Vì vậy chi phí mỗi tầng là \(\Oh(n\log n)\) theo
kích thước đoạn, và truy hồi tổng quát là
\[
  T(n)=2T(n/2)+\Oh(n\log n).
\]

\section{FJOI2012 Point}

Slide cuối nêu thêm một bài gợi ý. Có một mảng độ dài \(n\), ban đầu rỗng.
Thực hiện \(n\) thao tác \((P_i,Q_i)\), nghĩa là điền giá trị \(Q_i\) vào vị
trí \(P_i\). Hai dãy \(P_i\) và \(Q_i\) đều là hoán vị của \(1,\ldots,n\). Sau
mỗi thao tác, cần in số nghịch thế của dãy hiện tại; \(n\le 50000\).

Các cấu trúc như cây lồng cây hoặc danh sách khối có thể đạt gần
\(\Oh(n\log^2 n)\), nhưng trong kỳ thi chúng không dễ viết và không phải lúc
nào cũng đạt trọn điểm. Bài giảng dùng câu hỏi này để nhấn mạnh tinh thần của
CDQ: nếu tìm được một trục thời gian phù hợp, ta có thể hy vọng giảm bớt một
thừa số \(\log n\) bằng cách đưa các truy vấn động về xử lý offline.

\section{Ghi nhớ khi áp dụng}

Các ví dụ trong tài liệu có bề ngoài rất khác nhau, nhưng cùng dựa trên một
chuỗi quyết định:
\begin{enumerate}
  \item Xác định trục thời gian hoặc trục phụ thuộc: thao tác ở nửa trái có
  thể ảnh hưởng tới truy vấn hoặc trạng thái ở nửa phải.
  \item Đảm bảo ảnh hưởng nội bộ đã do đệ quy phụ trách.
  \item Biến bước ảnh hưởng chéo thành một bài toán offline: sắp xếp theo
  một chiều, quét, rồi dùng Fenwick, cây đoạn, bao lồi hoặc cấu trúc dữ liệu
  phù hợp cho các chiều còn lại.
  \item Phân tích truy hồi của bước giữa. Nếu bước giữa tốn
  \(\Oh(n\log^c n)\), thường nhận thêm một thừa số \(\log n\) ở tầng CDQ.
\end{enumerate}

\end{document}
